\section{Purpose}

A120 shows that the fivefold rotation $R_5$ of the icosahedron redistributes
the three $Z_3$ modes with the weights
\[
  p_0 = \tfrac29,\qquad p_1 = \tfrac{2+3\varphi}{9},\qquad p_2 = \tfrac{5-3\varphi}{9},
\]
and leaves one edge open: how these weights are connected with the mass
ratios of the charged leptons. This document closes that edge as far as the
corpus closes it today. It works out two routes, both of which give pure
ratios, without an anchor and without $K_\text{frak}$: the Koide formula from
a single matrix element of $R_5$, and the $\varphi$ skeleton of the weights
themselves. On the second route
\[
  \frac{m_\tau}{m_e} = \Bigl(74+\frac1{45}\Bigr)\varphi^8 = 3477.469
\]
agrees with the measured value to $3\times10^{-5}$, within the measurement
uncertainty. The document states just as clearly what the route does not
achieve: the factors $74$ and $1/45$ are observed, not derived, and for
$m_\mu/m_e$ there is no $\varphi$ route at this level.

\textbf{[STIPULATION]} As in A120, the embedding of the $C_3$ structure in the
icosahedral group $A_5$ is presupposed.

\section{Why the icosahedron: the three and the five}

\textbf{[B]} FFGFT is a $Z_3$ theory; the golden ratio is the number of the five.
The rotation group $A_5$ of the icosahedron is the smallest finite group in
which a threefold and a fivefold rotation occur together without commuting.
This non-commuting produces the mixing of the $Z_3$ modes under $R_5$ that A120
computes.

\textbf{[B]} The two symmetries enter at different levels. The three acts as a
lattice operation on the compact carrier: the triality of the $D_4$ lattice is
the $Z_3$ of the theory. The lattice knows no five, since
$|\mathrm{Aut}(D_4)| = 1152 = 2^7\cdot3^2$ and $5\nmid1152$; no icosahedron
acts on a $D_4$ lattice. The five enters only in the unrolled space
$\mathbb{R}^3\times S^1$ (A090), and it does so as a phase: the $72^\circ$
rotation has the characteristic polynomial
$(x-1)\bigl(x^2-(\varphi-1)x+1\bigr)$ with trace $\varphi$. Modulo~3,
$\varphi$ lies in $\mathrm{GF}(9)$ and is primitive there ($\varphi^4=-1$,
order~8); the eigenvalues $\zeta_5$ lie in $\mathrm{GF}(81)$, because
$5\mid80=|\mathrm{GF}(81)^*|$.

\textbf{[B]} The distribution over the two levels is forced, not chosen. In
characteristic~3 one has $x^3-1=(x-1)^3$; no $\mathrm{GF}(3^k)$ contains a
primitive third root of unity. The three therefore cannot be a phase and must
be lattice; the five cannot be lattice and must be phase. Three and five do
not meet in the geometry but in the field $\mathrm{GF}(81)$; the extension
$\mathrm{GF}(9)\to\mathrm{GF}(81)$ is the doubling $6=3\oplus3'$.

\section{One matrix element carries the Koide formula}

\textbf{[B]} The weight $p_0=2/9$ and the Koide angle $\theta=2/9$ (A110, A120)
are the same quotient $2/3^2$: in the numerator the number of non-trivial
$Z_3$ modes, in the denominator the square of the group order. For the matrix
element $A=\langle v_0\,|\,R_5\,v_e\rangle$ one has $|A|=\sqrt2/3$; the
modulus gives the Koide amplitude, the squared modulus the phase:
\[
  a_k = 1 + 3|A|\cos\bigl(|A|^2 + \tfrac{2\pi k}{3}\bigr),\qquad
  3|A| = \sqrt2,\qquad |A|^2 = \tfrac29 .
\]
The same number appears once as the modulus and once as the squared modulus of
the same matrix element. The Koide formula thus contains no free parameter any
more; the mass amplitudes follow from the $A_5$ geometry.

\textbf{[K]} The ratios $(a_0/a_1)^2$ and $(a_2/a_1)^2$ give
\begin{center}
\small
\begin{tabular}{@{}lllll@{}}
\toprule
Ratio & Formula & Value & Measured & Deviation \\
\midrule
$m_\tau/m_e$ & $(a_0/a_1)^2$ & $3477.473$ & $3477.37\pm0.18$ & $+3.1\times10^{-5}$ ($0.6\sigma$) \\
$m_\mu/m_e$ & $(a_2/a_1)^2$ & $206.77032$ & $206.7682827(46)$ & $+9.8\times10^{-6}$ ($442\sigma$) \\
\bottomrule
\end{tabular}
\end{center}
For $m_\tau/m_e$ the value lies within the measurement uncertainty of
$m_\tau$. For $m_\mu/m_e$ the exact pair $(\sqrt2,\,2/9)$ is excluded: at fixed
amplitude the ratio depends only on $\theta$, and the deviation exceeds the
uncertainty by more than two orders of magnitude.
This is not an objection to the origin of $2/9$ but a limit of the exact pair.

\textbf{[B]} Since $2/9\notin\pi\mathbb{Q}$, $\cos(2/9)$ is transcendental.
Exact algebraic mass amplitudes are therefore excluded in principle on this
route; any algebraic form can only approximate the Koide value.

\section{The $\varphi$ skeleton of the weights}

\textbf{[B]} The two weights containing $\varphi$ stand in exact ratios to each
other and to the rational weight:
\[
  \frac{p_1}{p_2} = \varphi^8,\qquad
  \frac{p_0}{p_2} = 2\varphi^4,\qquad
  \frac{p_1}{p_0} = \frac{\varphi^4}{2},\qquad
  3\sqrt{p_j}\in\bigl\{\sqrt2,\ \varphi^2,\ \varphi^{-2}\bigr\}.
\]
These are pure numbers from the geometry, without a cosine and without a
scale. The Koide amplitude $\sqrt2=3\sqrt{p_0}$ stands in the same row as
$\varphi^2$ and $\varphi^{-2}$.

\section{From the skeleton to the mass ratios}

\textbf{[K]} The mass ratios carry this skeleton. The most precise hit is
\[
  \frac{m_\tau}{m_e} = \Bigl(74+\frac1{45}\Bigr)\varphi^8 = \frac{3331}{45}\,\varphi^8 = 3477.469 ,
\]
that is $+3.0\times10^{-5}$ or $0.6\sigma$ against the measured value
$3477.37\pm0.18$ (PDG~2024). The overview:
\begin{center}
\small
\begin{tabular}{@{}llll@{}}
\toprule
Ratio & Formula & Value & against measured value \\
\midrule
$m_\tau/m_e$ & $74\,\varphi^8$ & $3476.42$ & $-2.7\times10^{-4}$ ($-5.3\sigma$) \\
$m_\tau/m_e$ & $(74+\tfrac1{45})\,\varphi^8$ & $3477.469$ & $+3.0\times10^{-5}$ ($0.6\sigma$) \\
$m_\tau/m_e$ & $74\,\varphi^8\,(1+\tfrac{27}{12}\xi)$ & $3477.468$ & $+2.9\times10^{-5}$ ($0.6\sigma$) \\
$m_\tau/m_e$ & Koide, $\theta=2/9$ & $3477.473$ & $+3.1\times10^{-5}$ ($0.6\sigma$) \\
$m_\mu/m_e$ & $30\,\varphi^4 = 15\,p_0/p_2$ & $205.62$ & $-0.55\,\%$ \\
\bottomrule
\end{tabular}
\end{center}

\textbf{[K]} On the $\varphi$ route, $m_\tau/m_e$ is therefore as precise as on
the Koide route, as a pure ratio in natural units. The route is algebraic: the
$\varphi$ form approximates the transcendental Koide value, and the second layer
($1/45$ or $\tfrac{27}{12}\xi$) carries the remainder. Against the Koide value
itself, $(74+\tfrac1{45})\varphi^8$ is off by only $1.2\times10^{-6}$.

\textbf{[S]} The factors are observed, not derived. $74=2\cdot37$ is the same
number as in the numerator of the rotation number $74/75=1-100\xi$ of the
recursion (A050) and in the value of $K_\text{frak}$ (A040). The second layer
has the same form as the Galois form $1/\alpha=3700/27=137+\tfrac1{27}$
(A130): an integer plus a unit fraction. The~37 appears only in numerator parts
($74=2\cdot37$, $3700=100\cdot37$), in no denominator ($45=3^2\cdot5$,
$27=3^3$); a common algebraic source is therefore not established.

\textbf{[B]} The~37 demonstrably does not come from the icosahedron. In the
phase channel it is the root of unity $\zeta_{37}$, which lives in
$\mathrm{GF}(3^{18})$, because the order of~3 modulo~37 is~18. The icosahedral
phase $\zeta_5$ lives in $\mathrm{GF}(81)=\mathrm{GF}(3^4)$, and $4\nmid18$:
$\mathrm{GF}(81)$ is not a subfield of $\mathrm{GF}(3^{18})$. The two fields
share only $\mathrm{GF}(9)$, that is, $\varphi$.

\section{Two readings of the same numbers}

\textbf{[B]} The ladder of A100 and the Koide route give the same numbers, but
they can be read in two directions. In reading~A the ladder
$r_k\,\xi^{p_k}$ is fundamental, and the step from the per-cent level to the
$10^{-5}$ level is a measured correction. In reading~B the Koide formula from
the icosahedron is fundamental; the ladder is then a rational approximation to
the transcendental amplitudes $a_k^2$, and the generation-dependent residuals
are the distance between cosine and power law, computed rather than measured
\textbf{[K]}. They differ for the three leptons because the three cosine values
differ. Which reading carries the weight is not decided by the calculation
\textbf{[S]}.

\section{What the route does not achieve}

\textbf{[S]} For $m_\mu/m_e$ there is no $\varphi$ route at the level of
$m_\tau/m_e$: $30\varphi^4$ is off by $0.55\,\%$, the ladder of A100 by
$0.52\,\%$, and the exact Koide pair is excluded (above).

\textbf{[K]} The $\varphi$ skeleton is not the earlier $\varphi$ mass ladder
with $m\propto\varphi^{\text{generation}}$ and $m_\mu/m_e\approx\varphi^{10}$.
That ladder is refuted: $\varphi^{10}=122.99$, not $206.77$.

\section{Verification Script}

\begin{itemize}
  \item Weights $p_j$ from $R_5$ (exactly against $2/9$, $(2+3\varphi)/9$,
    $(5-3\varphi)/9$), $|A|=\sqrt2/3$, the $\varphi$ skeleton, the Koide ratios
    against PDG~2024 and CODATA, the table of $\varphi$ forms, the field
    statements ($5\nmid1152$, $\varphi^4=-1$ in $\mathrm{GF}(9)$, order of~3
    modulo~37, $4\nmid18$) and the distinction from $\varphi^{10}$:
    \texttt{python/A\_Serie\_Skripte/a125\_phi\_skelett.py}
\end{itemize}

\section{Symbol Glossary}

Symbols used throughout the A series are explained in A010. Specific to this
document:

\begin{center}
\renewcommand{\arraystretch}{1.3}
{\small\begin{tabular}{lp{9cm}}
\toprule
Symbol & Meaning \\
\midrule
$\varphi$ & golden ratio $(1+\sqrt5)/2$ \\
$A_5$, $R_5$ & rotation group of the icosahedron; fivefold rotation by $72^\circ$ \\
$v_0$, $v_e$ & trivial and non-trivial (electron) mode of the $Z_3$ circulant \\
$p_j$ & redistribution weights $|\langle v_j\,|\,R_5\,v_e\rangle|^2$ \\
$A$ & matrix element $\langle v_0\,|\,R_5\,v_e\rangle$, $|A|=\sqrt2/3$ \\
$a_k$ & Koide amplitudes $1+\sqrt2\cos(\theta+2\pi k/3)$ \\
$\zeta_n$ & primitive $n$-th root of unity \\
\bottomrule
\end{tabular}}
\renewcommand{\arraystretch}{1.0}
\end{center}

\section{Open Questions}

\textbf{First, the factors $74$ and $1/45$.} They do not come from the weights
themselves and demonstrably not from the icosahedron. Whether they follow from
the Galois hierarchy $\mathrm{GF}(3^n)$, in which $K_\text{frak}$ and
$3700/27$ also stand, is open.

\textbf{Second, the common root.} The icosahedron supplies phases in
$\mathrm{GF}(81)$, number theory supplies orders in $\mathrm{GF}(3^{18})$; both
hit the lepton masses, but they touch only in $\varphi$. Whether there is a
structure in which both have the same root, or whether they are two
independent approaches, is the open edge.

\textbf{Third, $m_\mu/m_e$.} For this ratio an algebraic route at the level of
$10^{-5}$ is missing.

\textbf{Fourth, the embedding $C_3\subset A_5$.} As in A120, it is presupposed.

\section{Supersedes}

This document incorporates: Doc.~367 (one matrix element carries Koide),
Doc.~368 ($\varphi$ skeleton and Galois factors), Doc.~369 (routes~2 and~4 of
the inventory), Doc.~370 (why the icosahedron), Doc.~364 (as far as $A_5$ and
$D_4$ are concerned). Incorporated: R121. Delimited: Docs.~172, 281 (refuted
$\varphi$ mass ladder).

\section{Use of AI Tools}

This document was created with the assistance of AI tools. The questions,
stipulations, and all substantive decisions come from the author; the tools
formulate, compute, and create verification scripts.
Responsibility for content and errors rests entirely with the author.
Full wording of the rules in A010.
